\documentclass[a4paper,12pt]{article}
\usepackage[utf8]{inputenc}
\usepackage[T1]{fontenc}
\usepackage[french]{babel}
\usepackage{geometry}
\geometry{left=1.8cm,right=1.8cm,top=1.3cm,bottom=1.6cm}

\usepackage{lmodern}
\usepackage{xcolor}
\definecolor{titleblue}{RGB}{0,51,140}
\definecolor{TITLEBLUE}{RGB}{0,51,140}

\usepackage{hyperref}
\hypersetup{
    colorlinks=true,
    urlcolor=titleblue,
    linkcolor=titleblue,
    citecolor=titleblue,
    pdfborder={0 0 0}
}

\usepackage{titlesec}
\usepackage{enumitem}

\pagestyle{empty}

\titleformat{\section}
  {\color{titleblue}\Large\bfseries\uppercase}{}{0em}{}
  [\color{titleblue}\titlerule]

\titlespacing\section{0pt}{8pt}{4pt}
\setlist[itemize]{leftmargin=*,topsep=2pt,itemsep=1pt}

\begin{document}

\begin{center}
    {\Huge\bfseries\color{titleblue} AISSAOUI Ismail}\\[4pt]
    {\Large Ingénieur IA \& Data Science}\\[6pt]

    {\normalsize
    Email : \href{mailto:ismail.aissaoui.pro@gmail.com}{ismail.aissaoui.pro@gmail.com} \quad
    Tél : +213 660 70 77 96 \\
    Portfolio : \href{https://ismail-aissaoui.vercel.app}{ismail-aissaoui.vercel.app} \quad
    Localisation : Chlef, Algérie \\
    LinkedIn : \href{https://linkedin.com/in/aissaoui-ismail-6a77a92a7}{linkedin.com/in/aissaoui-ismail-6a77a92a7} \quad
    GitHub : \href{https://github.com/i-aissaoui}{github.com/i-aissaoui}
    }
\end{center}

%===========================
\section{RÉSUMÉ PROFESSIONNEL}
Ingénieur IA \& Data Science spécialisé dans les architectures LLM, le RAG et les systèmes de soutien cognitif. Passionné par l'interaction humain-IA et la conception de systèmes génératifs visant à étendre le raisonnement humain plutôt qu'à s'y substituer. Expert en déploiement de modèles de langage (Transformers) et en ingénierie de prompt avancée pour le soutien à la pensée critique.

%===========================
\section{COMPÉTENCES}
\setlength{\parskip}{-1pt}
{\normalsize
\textbf{Programmation}: Python, TypeScript, JavaScript, Java, SQL \\

\textbf{ML / DL}: PyTorch, TensorFlow, Keras, Scikit-learn, Lightning, \textbf{Fine-tuning (LoRA/QLoRA)} \\

\textbf{Techniques IA}: LLMs (GPT, Llama), RAG, \textbf{Multi-Agent Systems}, \textbf{Chain-of-Thought Reasoning}, GNN \\

\textbf{Optimisation}: PSO, Algorithme génétique, ACO, Recuit simulé, Optimisation bayésienne \\

\textbf{Web \& API}: React, Next.js, FastAPI, Node.js, TailwindCSS \\

\textbf{Data \& MLOps}: MLflow, Airflow, Kubeflow, DVC, Pandas, NumPy \\

\textbf{Bases de données}: PostgreSQL, MongoDB, MySQL, Neo4j, Cassandra \\

\textbf{Cloud \& DevOps}: Docker, Linux, Git, CI/CD, CUDA, AWS, Grafana \\

\textbf{Soft Skills}: Leadership technique, Communication, Gestion des parties prenantes
}

%===========================
\section{EXPÉRIENCE PROFESSIONNELLE}
\textbf{Stagiaire Ingénieur IA — Jumeau Numérique \& Maintenance Prédictive} \hfill 2026 \\
Sonatrach — Algérie \\
Développement d’un jumeau numérique pour turbines à gaz permettant une surveillance en temps réel et une maintenance prédictive. Implémentation d’une détection d’anomalies avancée et estimation de la RUL à l’aide de modèles de deep learning.

\textbf{Stagiaire en ingénierie réseau} \hfill Sep 2024 \\
Algérie Télécom RMC Center — Chlef, Algérie

%===========================
\section{FORMATION}
\textbf{École Nationale Supérieure d’Informatique (ESI-SBA)} \hfill 2021--2026 \\
Diplôme d’ingénieur et master en informatique \\
Spécialisation : Intelligence Artificielle \& Informatique

%===========================
\section{RÉALISATIONS CLÉS}
\begin{itemize}
    \item Développement de modèles Transformer (BERT, GPT, RAG) pour des applications multilingues réelles.
    \item Conception de plateformes d’IA explicable pour la HR tech et les institutions publiques.
    \item Système de recommandation basé sur graphes avec +24\% de précision (60\,000 profils).
    \item Pipelines ML complets : ingestion, validation, entraînement, CI/CD et monitoring.
    \item Vision temps réel (YOLO, ByteTrack) réduisant la latence de 32\%.
\end{itemize}

\newpage

%===========================
\section{PROJETS}
\setlength{\parskip}{13pt}
{\small

\textbf{Geo-RAG — Assistant Cognitif Cartographique} \hfill 2025 \\
Développement d'un système RAG visant à soutenir le raisonnement spatial et la découverte de connaissances géospatiales. Utilisation de l'IA pour augmenter l'analyse plutôt que d'automatiser la décision. \\
\textit{Tech : ChromaDB, LangChain, OpenAI — Concepts : RAG, Cognitive Support, Raisonnement spatial}

\textbf{Mini-GPT \& Architecture Transformer} \hfill 2025 \\
Implémentation d'un modèle GPT compact pour l'étude des mécanismes d'attention. Analyse de la manière dont les modèles génératifs encodent et restituent les structures de pensée. \\
\textit{Tech : PyTorch, CUDA, AdamW — Concepts : Transformers, Language Modeling, Latent Space}

\textbf{Pipeline NLP avancé multi-optimisation} \hfill 2025 \\
Modération multilingue optimisée via PSO, algorithme génétique et optimisation bayésienne. \\
\textit{Tech : DistilBERT, PyTorch, FastAPI — Concepts : Optimisation, Transformers, Classification}

\textbf{Planificateur universitaire intelligent} \hfill 2025 \\
Génération automatique d’emplois du temps sans conflits à l’aide d’heuristiques hybrides. \\
\textit{Tech : FastAPI, NumPy, Pandas — Concepts : GA, PSO, ACO, Recuit simulé}

\textbf{Mini-GPT (Decoder-only)} \hfill 2025 \\
Implémentation d’un GPT compact avec entraînement contrôlé et génération séquentielle. \\
\textit{Tech : PyTorch, CUDA, AdamW, W\&B — Concepts : Transformers, Language Modeling}

\textbf{North African Sentiment BERT} \hfill 2025 \\
Adaptation LoRA de BERT pour l’analyse de sentiment en dialecte maghrébin. \\
\textit{Tech : PyTorch, LoRA, MLflow — Concepts : Fine-tuning, Low-Rank Adaptation}

\textbf{Geo-RAG — Assistant cartographique} \hfill 2025 \\
Système RAG intégrant des données géospatiales pour répondre à des requêtes cartographiques. \\
\textit{Tech : ChromaDB, GeoPandas, Streamlit — Concepts : RAG, Raisonnement spatial}

\textbf{Recognizini — Reconnaissance faciale} \hfill 2025 \\
Reconnaissance faciale avec embeddings incrémentaux et boucle de feedback actif. \\
\textit{Tech : OpenCV, PyTorch, Redis, FastAPI — Concepts : Embeddings, Drift Management}

\textbf{Recognizini — IA Interactive \& Feedback} \hfill 2025 \\
Système de reconnaissance faciale avec boucle de feedback actif, plaçant l'humain au centre de l'apprentissage du modèle pour préserver l'autonomie et le contrôle. \\
\textit{Tech : OpenCV, PyTorch, FastAPI — Concepts : Human-in-the-loop, Agency, Active Learning}

\textbf{Suite analytique trafic} \hfill 2025 \\
Analyse du trafic routier : suivi, estimation de vitesse et OCR des plaques. \\
\textit{Tech : YOLOv8, ByteTrack, EasyOCR — Concepts : Tracking, OCR, Analytics}

\textbf{Système de gestion du Hadj} \hfill 2024 \\
Plateforme logistique end-to-end pour la gestion de groupes et d’itinéraires. \\
\textit{Tech : React, FastAPI, PostgreSQL — Concepts : Workflow, Synchronisation}

\textbf{POS \& e-commerce hybride} \hfill 2025 \\
Système POS offline-first intégré à une boutique en ligne. \\
\textit{Tech : PyQt, SQLite, PostgreSQL — Concepts : Sync offline/online, Inventory Management}
}

%===========================
\section{LANGUES}
Arabe : Langue maternelle \quad|\quad Anglais : Avancé \quad|\quad Français : Conversationnel

\end{document}
